\documentclass[11pt,a4paper]{article}
\usepackage{amsmath}
\usepackage{amsfonts}
\usepackage{amssymb}
\usepackage{graphicx}
\usepackage{float}

\title{Local Randomized Neural Networks with Discontinuous Galerkin Methods for Diffusive-Viscous Wave Equation}
\author{Summary of Sun and Wang (2023)}
\date{}

\begin{document}

\maketitle

\section{Introduction}

The diffusive-viscous wave equation (DVWE) extends traditional wave equation theory by incorporating both diffusion and viscosity effects. This equation plays a crucial role in geophysics applications, particularly in petroleum exploration, where it helps describe frequency-dependent reflection and fluid-saturation in porous media. Specifically, DVWE enables the detection of fluid-saturated layers by introducing frictional and viscous damping terms, making it essential for geological exploration.

Traditional numerical methods for solving DVWE include finite element methods and discontinuous Galerkin approaches. However, these methods often involve time discretization, which can lead to error accumulation over time steps. To address this limitation, this paper develops a novel approach that combines local randomized neural networks (LRNN) with discontinuous Galerkin (DG) methods to solve DVWE using a space-time framework.

The key advantages of this approach include:
\begin{itemize}
    \item Efficient handling of time-dependent problems without time discretization
    \item Avoidance of error accumulation over time iterations
    \item Higher accuracy with fewer degrees of freedom compared to traditional methods
    \item Effective solution of long-time simulation problems
\end{itemize}

\section{Neural Network Structure}

\subsection{Randomized Neural Networks}

The approach utilizes a special type of neural network called randomized neural networks (RNNs). The structure of these networks is defined as follows:

\begin{align}
U(x) &= W^{(L+1)}(N^{(L)} \circ \cdots N^{(2)} \circ N^{(1)}(x)) \\
N^{(i)}(x) &= \sigma(W^{(i)}x + b^{(i)}), \quad i = 1,2,\ldots,L
\end{align}

where:
\begin{itemize}
    \item $x \in \Omega$ is the input vector (spatial and temporal coordinates)
    \item $N^{(i)}$ is the $i$-th hidden layer with weight matrix $W^{(i)}$ and bias vector $b^{(i)}$
    \item $\sigma$ is a nonlinear activation function
    \item $L$ is the number of hidden layers (network depth)
    \item $W^{(L+1)}$ is the weight matrix of the output layer
\end{itemize}

The key innovation of RNNs is that, unlike standard neural networks where all parameters are trained through backpropagation, RNNs randomly assign and fix the parameters of the hidden layers. Only the weights of the output layer are trained, using a least-squares method. This approach significantly reduces training time while maintaining approximation capability.

The function space for RNNs can be written as:
\begin{equation}
M_{RNN}(D) = \left\{U(\alpha, \theta, x) = \sum_{j=1}^{M} \alpha_j^D \phi^D(\theta_j, x) : x \in D\right\}
\end{equation}

where:
\begin{itemize}
    \item $D \subset \Omega$ is the domain
    \item $\phi^D(\theta_j, x)$ is the output of the hidden layer with fixed parameters $\theta_j$
    \item $\alpha_j^D$ are the trainable weights of the output layer
    \item $M$ is the number of neurons in the last hidden layer
\end{itemize}

\subsection{Local Randomized Neural Networks}

For complex problems, a single neural network may not provide sufficient accuracy. Therefore, the domain is decomposed into multiple subdomains, with a separate RNN assigned to each subdomain. The outputs from the last hidden layer of each local network serve as basis functions, and the discontinuous Galerkin formulation is used to couple these local networks together.

\section{LRNN-DG Methods for Diffusive-Viscous Wave Equation}

\subsection{Problem Formulation}

The diffusive-viscous wave equation with mixed boundary conditions is given by:

\begin{align}
\frac{\partial^2 u}{\partial t^2} + \gamma \frac{\partial u}{\partial t} - \frac{\partial}{\partial t}\text{div}(\eta \nabla u) - \text{div}(\xi^2 \nabla u) &= f \quad \text{in } \Sigma \\
u &= g_D \quad \text{on } I \times \Gamma_D \\
\frac{\partial u}{\partial n} &= g_N \quad \text{on } I \times \Gamma_N \\
\frac{\partial u}{\partial n} + \kappa u &= g_R \quad \text{on } I \times \Gamma_R \\
u &= u_0 \quad \text{on } \{0\} \times \Omega \\
\frac{\partial u}{\partial t} &= w_0 \quad \text{on } \{0\} \times \Omega
\end{align}

where:
\begin{itemize}
    \item $\Omega \subset \mathbb{R}^d$ ($d = 1,2,3$) is the spatial domain
    \item $I = (0, T)$ is the temporal domain
    \item $\Sigma = I \times \Omega$ is the space-time domain
    \item $\partial\Omega = \Gamma_D \cup \Gamma_N \cup \Gamma_R$ is partitioned into Dirichlet, Neumann, and Robin boundaries
    \item $\gamma = \gamma(x) > 0$ is the diffusive attenuation coefficient
    \item $\eta = \eta(x) > 0$ is the viscous attenuation coefficient
    \item $\xi = \xi(x) > 0$ is the wave propagation speed
    \item $\kappa = \kappa(x) > 0$ is the coefficient for the Robin boundary condition
    \item $f(t,x)$ is the source term
\end{itemize}

\subsection{Domain Decomposition and Notation}

The space-time domain is partitioned as follows:
\begin{itemize}
    \item Time interval $I$ is divided into $N_t$ sub-intervals $D_\tau = \{I_i = (t_{i-1}, t_i), 0 = t_0 < t_1 < \cdots < t_{N_t} = T\}$
    \item Spatial domain $\Omega$ is decomposed into a mesh $\{T_h\}$ with $N_s$ elements
    \item The complete space-time domain $\Sigma$ is partitioned into $N_e = N_t \times N_s$ elements using $\{D_\tau \times T_h\}$
\end{itemize}

For spatial faces, we define averages $\{v\} = \frac{1}{2}(v^+ + v^-)$ and jumps $\llbracket v \rrbracket = v^+ n^+ + v^- n^-$ for scalar functions, and $\{q\} = \frac{1}{2}(q^+ + q^-)$ and $[q] = q^+ \cdot n^+ + q^- \cdot n^-$ for vector functions.

For temporal faces, we define averages $\{w(t_i,x)\} = \frac{1}{2}(w(t_i^+,x) + w(t_i^-,x))$ and jumps $[w(t_i,x)] = w(t_i^+,x) - w(t_i^-,x)$.

\subsection{LRNN-DG Method}

The LRNN-DG function spaces are defined as:
\begin{align}
V_h^\tau &= \{v_h \in L^2(\Sigma) : v_h^\tau|_\sigma \in M_{RNN}(\sigma) \text{ } \forall \sigma \in D_\tau \times T_h\} \\
Q_h^\tau &= \{q_h \in [L^2(\Sigma)]^d : q_h^\tau|_\sigma \in [M_{RNN}(\sigma)]^d \text{ } \forall \sigma \in D_\tau \times T_h\}
\end{align}

To derive the DG formulation, the DVWE is rewritten as a system of first-order equations:
\begin{align}
p &= \nabla u \quad \text{in } \Sigma \\
w &= \frac{\partial u}{\partial t} \quad \text{in } \Sigma \\
\frac{\partial w}{\partial t} + \gamma w - \nabla \cdot (\eta \nabla w) - \nabla \cdot (\xi^2 p) &= f \quad \text{in } \Sigma
\end{align}

Test functions are applied to these equations, integration by parts is performed, and numerical traces are introduced to approximate the solution across element boundaries. After combining the resulting equations and choosing appropriate numerical fluxes, the LRNN-DG scheme is formulated as: Find $u_h^\tau \in V_h^\tau$ such that
\begin{equation}
B_h^\tau(u_h^\tau, v_h^\tau) = l(v_h^\tau) \quad \forall v_h^\tau \in V_h^\tau
\end{equation}

where $B_h^\tau$ and $l$ are bilinear and linear forms containing terms for:
\begin{itemize}
    \item Temporal derivatives and initial conditions
    \item Diffusive terms with coefficient $\gamma$
    \item Viscous terms with coefficient $\eta$
    \item Wave propagation terms with coefficient $\xi^2$
    \item Interior penalty terms with parameters $\beta_1$ and $\beta_2$
    \item Boundary condition enforcement (Dirichlet, Neumann, and Robin)
\end{itemize}

This formulation results in a linear system $AU = L$, where $U$ represents the unknown parameters of the output layers of the local neural networks. This system is solved using a least-squares method.

\subsection{LRNN-C$^0$DG Method}

The LRNN-C$^0$DG method removes penalty terms by enforcing boundary conditions and continuity at selected collocation points:

\begin{enumerate}
    \item Dirichlet boundary condition at points $P_h^D = \{(t_D, x_D) \in D_\tau \times E_h^D\}$:
    \begin{equation}
    u_h^\tau(t_D, x_D) = g_D(t_D, x_D) \quad \forall (t_D, x_D) \in P_h^D
    \end{equation}
    
    \item Initial condition at points $P_\tau^I = \{(t_0, x_I) \in t_0 \times T_h\}$:
    \begin{equation}
    u_h^\tau(t_0, x_I) = u_0(x_I) \quad \forall (t_0, x_I) \in P_\tau^I
    \end{equation}
    
    \item C$^0$ continuity at spatial interior points $P_h^i = \{(t_i, x_i) \in (D_\tau \times E_h^i)\}$:
    \begin{equation}
    \llbracket u_h^\tau \rrbracket = 0 \quad \forall (t_i, x_i) \in P_h^i
    \end{equation}
    
    \item C$^0$ continuity at temporal interior points $P_\tau^i = \{(t_i, x_i) \in (P_\tau^i \times T_h)\}$:
    \begin{equation}
    [u_h^\tau(t_i, x_i)] = 0 \quad \forall (t_i, x_i) \in P_\tau^i
    \end{equation}
\end{enumerate}

These constraints form a system of equations $A_2 U = L_2$. Combined with the weak formulation equations $A_1 U = L_1$, the complete system is solved using least squares:
\begin{equation}
\begin{bmatrix} A_1 \\ A_2 \end{bmatrix} U_0 = \begin{bmatrix} L_1 \\ L_2 \end{bmatrix}
\end{equation}

\subsection{LRNN-C$^1$DG Method}

The LRNN-C$^1$DG method extends the C$^0$ continuity to C$^1$ continuity by imposing additional constraints on the gradient and time derivative at collocation points:

\begin{enumerate}
    \item Spatial gradient continuity:
    \begin{equation}
    [\nabla_h u_h^\tau] = 0 \quad \forall (t_i, x_i) \in P_h^i
    \end{equation}
    
    \item Temporal derivative continuity:
    \begin{equation}
    \left[\frac{\partial^\tau u_h^\tau}{\partial^\tau t}(t_i, x_i)\right] = 0 \quad \forall (t_i, x_i) \in P_\tau^i
    \end{equation}
\end{enumerate}

The complete LRNN-C$^1$DG system combines these additional constraints with the weak formulation and the C$^0$ constraints, resulting in another least-squares problem.

\section{Theoretical Findings}

The paper provides a consistency proof for the LRNN-DG method:

\begin{theorem}[Consistency]
The LRNN-DG scheme is consistent, i.e., for the solution $u \in C^1(I; H^2(\Omega))$ of the DVWE, we have
\begin{equation}
B_h^\tau(u, v_h^\tau) = l(v_h^\tau) \quad \forall v_h^\tau \in V_h^\tau
\end{equation}
\end{theorem}

The proof uses integration by parts and the properties of the exact solution to show that the bilinear form applied to the exact solution equals the linear form for any test function.

For the LRNN-C$^0$DG and LRNN-C$^1$DG methods, a numerical experiment in a previous paper demonstrates that the jump terms $\llbracket u_h^\tau \rrbracket$, $[u_h^\tau]$, and boundary mismatch $u_h^\tau - g_D$, $u_h^\tau - u_0$ approach zero in the $L^2$ norm when sufficient collocation points are used.

\section{Numerical Examples and Results}

\subsection{One-Dimensional Example}

The one-dimensional diffusive-viscous wave equation was solved on domain $\Omega = (0,1)$ with time interval $I = (0,1)$ and Dirichlet boundary conditions. Coefficients were chosen to represent water-saturated rock: $\gamma = 90$, $\eta = 2 \times 10^{-7}$, and $\xi = 1.47$. The exact solution was $u(t,x) = e \cos(27x^2 + 27t^2 + 8\pi - 24)(x^2 + 1)^2(t^2 + 1)^2$.

Results showed:
\begin{itemize}
    \item All three methods (LRNN-DG, LRNN-C$^0$DG, LRNN-C$^1$DG) achieved high accuracy
    \item With 320 degrees of freedom per subdomain and mesh size $h = 1/10$, relative $L^2$ errors reached $6.57 \times 10^{-5}$ to $7.53 \times 10^{-5}$
    \item Errors decreased consistently as mesh size decreased and degrees of freedom increased
    \item Absolute error distributions showed the methods effectively captured oscillatory solutions
\end{itemize}

\subsection{Two-Dimensional Example}

A two-dimensional diffusive-viscous wave equation was solved on domain $\Omega = (0,1)^2$ with time interval $I = (0,0.5)$ and Dirichlet boundary conditions. Parameters were set to $\gamma = 1$, $\eta = 0.01$, and $\xi = 0.1$. The exact solution was $u(t,x,y) = e^{-t} \sin(2\pi x) \sin(2\pi y)$.

Results demonstrated:
\begin{itemize}
    \item With 320 degrees of freedom per element and mesh size $h = 1/4$, relative $L^2$ errors reached $6.82 \times 10^{-5}$ to $1.20 \times 10^{-5}$
    \item Fixed time sub-interval $\tau = 1/4$ with smaller spatial mesh sizes achieved similar accuracy with fewer total degrees of freedom
    \item The methods were competitive with traditional local DG methods but required only a single linear least-squares solve rather than many time iterations
\end{itemize}

\subsection{Three-Dimensional Example}

A three-dimensional diffusive-viscous wave equation was solved on domain $\Omega = (0,1)^3$ with time interval $I = (0,5)$ and mixed boundary conditions. Parameters for dry sandstone were used: $\gamma = 56$, $\eta = 5.6 \times 10^{-8}$, and $\xi = 1.19$. The exact solution was $u(t,x,y,z) = t^2 \sin(\pi x) \sin(\pi y) \sin(\pi z)$.

A direct comparison with the interior penalty DG and Crank-Nicholson (IPDG-CN) method showed:
\begin{itemize}
    \item LRNN-DG achieved a relative $L^2$ error of $2.22 \times 10^{-4}$ in 33.5 seconds
    \item IPDG-CN achieved a similar relative $L^2$ error of $2.80 \times 10^{-4}$ but required 15,400 seconds
    \item LRNN-DG outperformed IPDG-CN in both accuracy and computational efficiency by several orders of magnitude
\end{itemize}

\section{Implications and Advantages}

\subsection{Computational Efficiency}

\begin{itemize}
    \item \textbf{Space-time approach}: The methods treat temporal and spatial variables equally, avoiding time-stepping procedures and error accumulation
    \item \textbf{Single-solve advantage}: Only one linear least-squares problem needs to be solved, regardless of the simulation time length
    \item \textbf{Training simplicity}: Only output layer weights are trained, drastically reducing the number of parameters compared to standard neural networks
    \item \textbf{Performance scaling}: Computational advantages become more pronounced for long-time simulations
\end{itemize}

\subsection{Accuracy and Solution Quality}

\begin{itemize}
    \item \textbf{High accuracy per DoF}: The methods achieve better accuracy than traditional approaches with the same degrees of freedom
    \item \textbf{Error reduction strategies}: Accuracy can be improved by either increasing degrees of freedom per subdomain or decreasing mesh size
    \item \textbf{Collocation effectiveness}: Enforcing continuity through collocation points proves to be highly effective
    \item \textbf{Complex solution handling}: The methods effectively capture oscillatory solutions and steep gradients
\end{itemize}

\subsection{Applicability to Geophysical Problems}

\begin{itemize}
    \item \textbf{Material property handling}: Successfully models different materials (water-saturated rock, dry sandstone) by adjusting coefficients
    \item \textbf{Mixed boundary conditions}: Naturally accommodates Dirichlet, Neumann, and Robin boundary conditions
    \item \textbf{Long-time simulations}: Particularly advantageous for geophysical applications requiring extended time horizons
    \item \textbf{Dimensional flexibility}: Methods scale effectively from 1D to 3D problems
\end{itemize}

\subsection{Challenges and Future Directions}

\begin{itemize}
    \item \textbf{Activation function selection}: Determining optimal activation functions for different problem types
    \item \textbf{Hidden layer parameter optimization}: Moving beyond uniform random distributions for parameter initialization
    \item \textbf{Model depth exploration}: Investigating deeper network architectures for improved approximation capability
    \item \textbf{Parallelization potential}: Exploiting the local nature of the approach for parallel computing
\end{itemize}

\end{document} 